\chapter{从单轨迹到实验：统计平均、可观测量与适用范围}
\label{chap:observables}

到目前为止，$P_n$ 都对应给定轨迹、给定到达时刻和给定初态。实验却可能同时包含电子横向束斑、到达时间分布、有限收集孔径以及量子相干。最后一步因此不是再改 Hamiltonian，而是明确“哪些自由度在振幅层相干求和，哪些自由度已经成为经典统计混合，哪些自由度被探测器迹掉”。只有先回答这一问题，才能合法地对 $J_n^2(2|\beta|)$ 做空间或时间平均。

\section{纯态、统计混合与探测器积分的区别}

在写具体圆形孔径积分之前，先区分三种常被混在一起的平均：单个纯电子波包的量子平均、实验束流的统计混合、探测器孔径积分。对给定轨迹标签 $\lambda$（可同时包含横向冲量参数、到达时刻、初始能量等），一阶 PINEM 演化给出边带振幅
\begin{equation}
 c_n(\lambda)
 =\ee^{\ii\vartheta_n(\lambda)}
 J_n\!\left(2|\beta(\lambda)|\right),
 \label{eq:general-trajectory-sideband-amplitude}
\end{equation}
其中相位 $\vartheta_n$ 由 $\arg\beta$ 和载波参考相位决定，而概率只含
$J_n^2(2|\beta|)$。

若入射电子是统计混合
\begin{equation}
 \widehat\rho_{\rm e}^{\rm in}
 =\int\dd\lambda\,W(\lambda)
 |\lambda\rangle\langle\lambda|,
 \qquad
 W(\lambda)\ge0,
 \qquad
 \int\dd\lambda\,W(\lambda)=1,
 \label{eq:electron-incoherent-mixture}
\end{equation}
且探测器不分辨 $\lambda$，则 Born 规则直接给出
\begin{equation}
 \boxed{
 P_n^{\rm mix}
 =\int\dd\lambda\,W(\lambda)
 J_n^2\!\left(2|\beta(\lambda)|\right).}
 \label{eq:general-incoherent-average}
\end{equation}
这就是后面所有“对电子束位置、延迟或孔径作概率平均”的一般形式。

但若 $\lambda$ 是同一个纯态内部保持相干的连续自由度，就不能先取模平方。设
\begin{equation}
 |\psi_{\rm in}\rangle
 =\int\dd\lambda\,g(\lambda)|\lambda\rangle,
 \qquad
 \int\dd\lambda\,|g(\lambda)|^2=1.
\end{equation}
第 $n$ 阶边带的未归一化态为
\begin{equation}
 |\psi_n\rangle
 =\int\dd\lambda\,g(\lambda)c_n(\lambda)
 |\lambda;n\rangle,
\end{equation}
故一般概率是双积分
\begin{equation}
 \boxed{
 P_n
 =\int\dd\lambda\dd\lambda'\,
 g(\lambda)g^*(\lambda')
 c_n(\lambda)c_n^*(\lambda')
 \mathcal K_n(\lambda',\lambda),}
 \label{eq:coherent-average-kernel}
\end{equation}
其中
\begin{equation}
 \mathcal K_n(\lambda',\lambda)
 \equiv\langle\lambda';n|\lambda;n\rangle
\end{equation}
是探测后剩余自由度的相干核。只有当不同 $\lambda$ 的最终态近似正交，
\begin{equation}
 \mathcal K_n(\lambda',\lambda)
 \simeq\delta(\lambda-\lambda'),
 \label{eq:diagonal-detection-limit}
\end{equation}
式~\eqref{eq:coherent-average-kernel} 才退化为式~\eqref{eq:general-incoherent-average}。因此“把 $J_n^2$ 在波包上积分”隐含的是冻结轨迹/正交探测或统计混合条件；若研究电子纵向相干、Talbot 演化、两次 PINEM 相位干涉或横向模式相干，就应保留双点核而不是直接做概率平均。

\section{从一般平均到实验能谱与 PINEM 图像}

实验能谱通常属于式~\eqref{eq:general-incoherent-average} 的情形：电子束在横向冲量参数上可用经典束流分布 $W_\perp(\bm R_\perp)$ 描述，探测器只记录能量。最一般的孔径平均因此是
\begin{equation}
 \boxed{
 P_n^{\rm det}
 =\frac{\displaystyle
 \int_{\mathcal A}\dd^2R_\perp\,
 W_\perp(\bm R_\perp)
 \int\dd t_c\,W_t(t_c)
 J_n^2\!\left(2|\beta(\bm R_\perp,t_c)|\right)}
 {\displaystyle
 \int_{\mathcal A}\dd^2R_\perp\,W_\perp(\bm R_\perp)}.}
 \label{eq:general-detector-average}
\end{equation}
下面的圆形窗口、均匀横向权重和 Gaussian 延迟分布都是这一母式的特例。

实验能谱不是固定 \((b,\varphi)\) 的单轨迹概率，而是孔径窗口内所有轨迹的平均。对半径 \(a\) 的球形投影和半径 \(w\) 的圆形采集窗口，归一化面积为
\begin{equation}
A_w=\pi(w^2-a^2),
\end{equation}
故
\begin{equation}
P_{\rm EEGS}(n)
=\frac{1}{\pi(w^2-a^2)}
\int_0^{2\pi}\dd\varphi\int_a^w b\dd b\,P_n(b,\varphi).
\label{eq:EEGS_def}
\end{equation}
将单轨迹概率式~\eqref{eq:Pn-gaussian} 代入空间平均定义，而不作级数展开，得到可直接数值计算的形式
\begin{equation}
\boxed{
\begin{aligned}
P_{\rm EEGS}(n;\tau)
={}&\frac{1}{\pi(w^2-a^2)}
\int_0^{2\pi}\dd\varphi\int_a^w b\dd b
\int_{-\infty}^{\infty}\dd z'\\
&\times
\left|g(z',-\infty)
J_n\!\left[
-\frac{e}{\hbar\omega}
|\Ftilde(b,\varphi)|
\exp\!\left(-\frac{(z'+v_0\tau)^2}
{4v_0^2\sigma_{\mathrm L}^2}\right)
\right]\right|^2.
\end{aligned}}
\label{eq:detector-averaged-spectrum}
\end{equation}
这个三重积分是最直接的实验能谱公式；下面的 \(D_{jk}^n\) 级数只是把其中的 \(b,\varphi\) 依赖在指数近似下进一步解析积分掉。

若需要把式~\eqref{eq:detector-averaged-spectrum} 在指数型横向近场模型下继续化成 Gamma/Beta 函数与双重级数的闭式，必须先建立 Bessel--Taylor 系数 $C_j^n$。为保持正文单向闭合，这一步不在这里提前调用尚未定义的级数，而统一放到附录的解析级数推导中。直接数值计算时，式~\eqref{eq:detector-averaged-spectrum} 已经是完整的可观测量公式。

对 PINEM 图像，固定 \((b,\varphi)\) 而把所有增益边带相加。由于归一化
\(
\sum_{n=-\infty}^{\infty}P_n=1
\)
且 \(P_n=P_{-n}\)，有
\begin{align}
1
&=P_0+2\sum_{n=1}^{\infty}P_n,\\
\boxed{
I_{\rm PINEM}(b,\varphi)
=\sum_{n=1}^{\infty}P_n(b,\varphi)
=\frac12[1-P_0(b,\varphi)].}
\label{eq:PINEMimage}
\end{align}
因此空间图像、偏振图像与 EEGS/EELS 能谱都来自同一组 \(P_n\)，差别只在于对空间、偏振和边带指标采用不同的求和或积分。



\section{近似层级、失效路径与实际计算顺序}

主线已经闭合。下面不再引入新理论，只把前文局部出现的近似条件放回它们各自的母方程中。判断某个实验参数能否使用下游公式时，应从上到下检查；一旦某一行失效，就回到该行左侧的上一级方程，而不是继续套用后面的特殊结果。

\begin{center}
\small
\begin{tabularx}{\textwidth}{>{\raggedright\arraybackslash}p{0.14\textwidth}>{\raggedright\arraybackslash}X>{\raggedright\arraybackslash}p{0.22\textwidth}>{\raggedright\arraybackslash}X}
\hline
近似步骤 & 母方程 & 控制量 & 失效后恢复什么\\
\hline
正能单带化
& Feshbach 精确式~\eqref{eq:Feshbach-exact-Heff-time}
& $\|H_{-+}\|/(2\mathcal E_0)$、$\hbar\Omega_{\rm drv}/(2\mathcal E_0)$
& 负能分量与 resolvent 下一阶\\
窄带包络
& 色散 Taylor 式~\eqref{eq:E-taylor-vector}
& $\Delta p/p_0$
& 三阶色散、基底随动修正\\
固定直线轨迹
& 三维包络式~\eqref{eq:rel-envelope-tensor}
& $\epsilon_{\perp,{\rm diff}},\epsilon_{\perp,{\rm kick}}$
& 横向衍射与轨迹弯折\\
一阶 eikonal
& 二阶包络式~\eqref{eq:rel-master}
& $\epsilon_{\rm disp},\epsilon_{\rm grad},\epsilon_{A^2},\epsilon_{\nabla A}$
& 色散、场梯度、$A^2$\\
无反冲通道
& 精确式~\eqref{eq:exact-channel-momentum}
& $\epsilon_{{\rm rec},n}$
& 非线性 $p_n(\mathcal E_n)$\\
准单色载波
& 时域积分式~\eqref{eq:g-integral}
& $\epsilon_{\rm SVEA}^{\max}$
& 包络导数对 $A$--$E$ 的修正\\
冻结空间场型
& 全频谱 Maxwell 解
& $\epsilon_{\rm mode}$
& 保留 $\bm E_{\Omega}(\bm r)$ 的频率依赖\\
冻结光包络
& 式~\eqref{eq:pulse-freeze-exact}
& $T_{\rm tr}/\sigma_{\rm L}$
& 近场区内部的包络变化\\
边带分离
& 振幅级 Fourier 变换
& $\Delta k_{\rm env}/(\omega/v_0)$
& 相邻边带振幅干涉\\
Rayleigh
& 精确 Mie/$T$ 算符
& $|k_{\rm h}|a$
& 高阶多极、迟滞与辐射修正\\
\hline
\end{tabularx}
\end{center}

这些条件属于不同层次，不能相互替代。例如 $|k_{\rm h}|a\ll1$ 只决定电磁散射能否约化成 Rayleigh 偶极，与电子是否无反冲没有逻辑关系；$|\beta|\ll1$ 只决定 Bessel 分布能否截到单光子阶，也不是推导一阶 eikonal 方程的必要条件。只要式~\eqref{eq:first-order-allconditions} 成立，即使 $|\beta|$ 不小，仍应保留式~\eqref{eq:Pn-beta-standard} 的完整 Bessel 级数，而不是把“强边带”误解成“电子动力学必然超出一阶”。


\subsection{一套不会逆序的实际计算流程}
对于一个新的 PINEM 几何，可以按下面的顺序工作：
\begin{enumerate}
\item 先给定电子中心能量、$\bm v_0$、波包宽度和光学频率；由第~\ref{chap:electron-reduction} 章检查正能单带化和一阶 eikonal 条件。
\item 独立求解 Maxwell 边值问题。解析球使用第~\ref{chap:mie-sphere} 章；一般几何使用第~\ref{chap:maxwell-general} 章的 Green/$T$ 算符或数值 Maxwell 求解器。
\item 把求得的复频域场逐点代入式~\eqref{eq:Fparallel-general-def}，得到 $\mathcal F_{\parallel}$ 和 $\beta$。这一步才把电磁模块接到电子模块。
\item 单色相干条件下用式~\eqref{eq:Pn-beta-standard} 计算边带；需要反冲时使用真实 $p_n(\mathcal E_n)$，而不是线性动量梯。
\item 有限脉冲时从第~\ref{chap:pulses} 章的局域 $\beta$ 开始；实验束流和探测器平均最后按第~\ref{chap:observables} 章处理。
\end{enumerate}
任何步骤都不应倒过来。例如，不能先假定 Rayleigh 极化率再声称得到了“一般 Mie”；也不能先把电子定义成等间隔动量梯再引入反冲；更不能把探测器概率平均塞回单电子 Schrödinger/Dirac 动力学中。

